\section{Results}
\label{sec:results}

\subsection{Explanation quality}
\label{sec:quality}

Table~\ref{tab:e1} reports lexical and embedding overlap between generated and reference explanations.
Zero-shot systems cluster tightly (BLEU-4 $0.016$--$0.026$, BERTScore $0.418$--$0.547$); the best of them, LLaVA-NeXT-Video-7B, is an order of magnitude below the adapted model on BLEU-4.
Size does not explain the ordering: Qwen2.5-VL-7B trails its 3B sibling on every lexical metric because it spends its 256-token budget on markdown headers before reaching any narrative.

\begin{table*}[!t]
\centering
\caption{Explanation quality on the CCD test split ($n{=}150$), computed on the Explanation field (whole output for systems that emit no field). BLEU-4 is corpus-level; other metrics are per-video means. $\uparrow$: higher is better. Best per column in \textbf{bold}. No \tcd{} variant differs significantly from greedy on any column (Table~\ref{tab:paired}).}
\label{tab:e1}
\small
\setlength{\tabcolsep}{8pt}
\begin{tabular*}{\textwidth}{@{\extracolsep{\fill}}lccccc@{}}
\toprule
\textbf{System} & \textbf{BLEU-4}$\uparrow$ & \textbf{ROUGE-L}$\uparrow$ & \textbf{METEOR}$\uparrow$ & \textbf{CIDEr}$\uparrow$ & \textbf{BERTScore}$\uparrow$ \\
\midrule
ZS Qwen2.5-VL-7B        & 0.016 & 0.160 & 0.203 & 0.338 & 0.486 \\
ZS Qwen2.5-VL-3B        & 0.019 & 0.180 & 0.234 & 0.394 & 0.518 \\
ZS Qwen2-VL-2B          & 0.018 & 0.120 & 0.132 & 0.221 & 0.418 \\
ZS LLaVA-NeXT-Video-7B  & 0.026 & 0.197 & 0.275 & 0.449 & 0.547 \\
\midrule
\crashlogic{} (greedy)                  & \textbf{0.142} & 0.336 & \textbf{0.371} & 0.508 & \textbf{0.686} \\
\crashlogic{} + \tcd{} ($\alpha{=}0.5$) & \textbf{0.142} & \textbf{0.339} & 0.369 & \textbf{0.515} & \textbf{0.686} \\
\crashlogic{} + \tcd{} ($\alpha{=}1.0$) & 0.140 & 0.338 & 0.365 & 0.507 & 0.684 \\
\bottomrule
\end{tabular*}
\end{table*}


\crashlogic{} reaches BLEU-4 $0.142$, ROUGE-L $0.336$, METEOR $0.371$, CIDEr $0.508$, BERTScore $0.686$.
Paired against zero-shot Qwen2.5-VL-7B, per-video BLEU improves on 149 of 150 videos and ROUGE-L on 150 of 150 ($p<10^{-4}$; Table~\ref{tab:paired}).
Adding \tcd{} at $\alpha{=}0.5$ leaves BLEU-4 and BERTScore unchanged to three decimals and moves ROUGE-L by $+0.003$ (paired 95\% CI $[-0.005,+0.011]$).
Fine-tuning is responsible for all lexical gains; no decoding variant changes any lexical metric significantly in the positive direction (Section~\ref{sec:ablation}).

\subsection{Entailment against the forensic ground truth}
\label{sec:faith}

We also score each output with a DeBERTa-v3 NLI cross-encoder (full matrix in Table~\ref{tab:e2}).
Two readings matter.
With the extracted Explanation as hypothesis, fine-tuning raises entailment mass ($P_e$: $0.139\to0.224$, $p{=}0.007$) but raises contradiction mass by more ($P_c$: $0.092\to0.229$, $p<10^{-4}$), so NLI-Score alone is flat or slightly worse (Table~\ref{tab:paired}).
Zero-shot models mostly say nothing and are scored as neutral rather than wrong.
With the \emph{entire} output as hypothesis (Full-NLI), the ordering reverses: \crashlogic{} greedy scores $0.247$ against $0.005$--$0.111$ for zero-shot ($\Delta{=}{+}0.170$ vs.\ Qwen2.5-VL-7B, $p{=}0.0005$).
We therefore treat NLI as a secondary signal and place the primary faithfulness claims on the structured costs below.

\subsection{Omission versus hallucination}
\label{sec:ho}

Table~\ref{tab:ho} is the central result.
It separates what each system fails to say from what it says wrongly, with bootstrap intervals on both.

\begin{table*}[!t]
\centering
\caption{Omission and hallucination costs (Eq.~\eqref{eq:ho}; $\downarrow$ better) with 95\% bootstrap intervals, $n{=}150$. \emph{Assert.} is the mean fraction of the five forensic fields a system asserts; \Hcost{} is conditional on assertion. Paired tests vs.\ zero-shot Qwen2.5-VL-7B (adaptation) and vs.\ greedy (decoding) are in Table~\ref{tab:paired}.}
\label{tab:ho}
\small
\setlength{\tabcolsep}{8pt}
\begin{tabular*}{\textwidth}{@{\extracolsep{\fill}}lccc@{}}
\toprule
\textbf{System} & \textbf{Assert.} & \textbf{\Ocost{}}$\downarrow$ & \textbf{\Hcost{}}$\downarrow$ \\
\midrule
ZS Qwen2.5-VL-7B        & 0.61 & 0.462 $[.436,.488]$ & 0.227 $[.188,.268]$ \\
ZS Qwen2.5-VL-3B        & 0.68 & 0.413 $[.391,.436]$ & 0.285 $[.247,.324]$ \\
ZS Qwen2-VL-2B          & 0.33 & 0.607 $[.570,.642]$ & 0.564 $[.495,.630]$ \\
ZS LLaVA-NeXT-Video-7B  & 0.86 & 0.334 $[.314,.354]$ & 0.401 $[.371,.432]$ \\
\midrule
\crashlogic{} (greedy)                  & 1.00 & 0.107 $[.088,.127]$ & 0.223 $[.200,.248]$ \\
\crashlogic{} + \tcd{} $\alpha{=}0.5$   & 1.00 & 0.102 $[.082,.122]$ & 0.212 $[.185,.236]$ \\
\crashlogic{} + \tcd{} $\alpha{=}1.0$   & 1.00 & 0.112 $[.093,.132]$ & 0.227 $[.201,.255]$ \\
\crashlogic{} + \tcd{} $\alpha{=}1.5$   & 1.00 & 0.124 $[.103,.148]$ & 0.249 $[.220,.276]$ \\
\crashlogic{} + \tcd{} $\alpha{=}2.0$   & 1.00 & 0.128 $[.106,.150]$ & 0.253 $[.225,.282]$ \\
\crashlogic{} + \tcd{} shuffle          & 1.00 & 0.114 $[.093,.136]$ & 0.233 $[.205,.262]$ \\
\crashlogic{} + SEASON (full)           & 1.00 & 0.119 $[.099,.140]$ & 0.240 $[.213,.267]$ \\
\bottomrule
\end{tabular*}
\end{table*}


\paragraph{Adaptation eliminates omission.}
Zero-shot Qwen2.5-VL-7B omits $46\%$ of the forensic content the ground truth specifies; \crashlogic{} omits $11\%$.
The paired difference is $-0.356$ (95\% CI $[-0.382,-0.328]$, $p<10^{-4}$), and the omission cost falls on 140 of 150 videos and rises on none.
Per field (Table~\ref{tab:fields-full}): the window is missing or badly placed in $98\%$ of zero-shot outputs against $43\%$ after adaptation; weather in $84\%$ against $12\%$; vehicles in $19\%$ against $2\%$.
The serialised target format does what it was designed to do.

\paragraph{Adaptation does not reduce hallucination.}
The hallucination cost of \crashlogic{} greedy is $0.223$ (95\% CI $[0.200,0.248]$).
Zero-shot Qwen2.5-VL-7B scores $0.227$ ($[0.188,0.268]$).
The paired difference is $-0.005$ ($[-0.051,+0.041]$, $p{=}0.81$; 67 videos better, 68 worse).
The two systems assert very different amounts---the adapted model emits all five fields on every video, the zero-shot model $61\%$ of them on average---but the fraction of asserted claims that contradict the ground truth is the same.
Domain adaptation has converted silence into assertions, roughly one in five of which is false.
The finding echoes Gekhman~\emph{et al.}~\cite{gekhman2024finetuning} for text-only LLMs: supervised targets the model cannot ground are learned as a policy of confident answering.
LLaVA-NeXT-Video-7B and the smaller Qwen models hallucinate more than either ($0.285$--$0.564$), so adaptation \emph{is} the best available system on \Hcost{}; the point is that it does not improve on its own backbone.

\subsection{Temporal grounding}
\label{sec:temporal}

Table~\ref{tab:e3} reports crash-window metrics with the prior-only baselines of Section~\ref{sec:metrics-temporal}.
Zero-shot systems emit a parseable window in $9$--$51\%$ of videos and score tIoU $\leq0.024$.
\crashlogic{} emits one on every video and scores $0.373$ (THR@0.50 $=0.440$); \tcd{} at $\alpha{=}0.5$ scores $0.394$ (THR@0.50 $=0.473$).

\begin{table*}[!t]
\centering
\caption{Crash-window grounding on the CCD test split ($n{=}150$; the 9 no-crash videos score tIoU $=0$ for every row). 95\% bootstrap CI on tIoU in brackets. \crashlogic{} emits only two windows, $[2,3]$\,s (88 videos) and $[3,4]$\,s (53); predicted vs.\ true start Spearman $\rho{=}0.009$ ($p{=}0.91$). Zero-shot parse rates are $0.09$--$0.51$. Constant-window baselines ignore the video.}
\label{tab:e3}
\small
\setlength{\tabcolsep}{8pt}
\begin{tabular*}{\textwidth}{@{\extracolsep{\fill}}lccc@{}}
\toprule
\textbf{System} & \textbf{tIoU}$\uparrow$ & \textbf{THR@.25}$\uparrow$ & \textbf{THR@.50}$\uparrow$ \\
\midrule
ZS Qwen2.5-VL-7B        & 0.012 $[.000,.028]$ & 0.000 & 0.000 \\
ZS Qwen2.5-VL-3B        & 0.010 $[.003,.018]$ & 0.000 & 0.000 \\
ZS Qwen2-VL-2B          & 0.000 $[.000,.000]$ & 0.000 & 0.000 \\
ZS LLaVA-NeXT-Video-7B  & 0.024 $[.007,.045]$ & 0.000 & 0.000 \\
\midrule
\crashlogic{} (greedy)                  & 0.373 $[.312,.438]$ & 0.573 & 0.440 \\
\crashlogic{} + \tcd{} ($\alpha{=}0.5$) & 0.394 $[.336,.456]$ & 0.600 & 0.473 \\
\midrule
Constant $[2,3]$\,s                     & 0.366 & 0.567 & 0.433 \\
Constant $[3,4]$\,s (training mode)     & \textbf{0.408} & 0.593 & 0.473 \\
Constant $[2,4]$\,s                     & \textbf{0.518} & \textbf{0.873} & \textbf{0.767} \\
\bottomrule
\end{tabular*}
\end{table*}


The baselines reframe these numbers.
The constant window $[3,4]$\,s---the most frequent training window---scores tIoU $0.408$ and THR@0.50 $=0.473$, matching \tcd{} exactly on the hit rate and exceeding both adapted systems on mean tIoU.
The constant window $[2,4]$\,s scores $0.518$.
\crashlogic{} emits only two distinct windows, $[2,3]$\,s (88 crash videos) and $[3,4]$\,s (53), and which one it picks is unrelated to when the crash happens (Spearman $\rho{=}0.009$ between predicted and true start, $p{=}0.91$).
Section~\ref{sec:prior} examines this further.
The honest summary of Table~\ref{tab:e3} is that adaptation teaches the model the CCD temporal marginal, and that tIoU without a prior baseline would have hidden this.

\subsection{Does temporal contrastive decoding help?}
\label{sec:tcd-results}

Table~\ref{tab:paired} gives paired statistics for the key comparisons.
At $\alpha{=}0.5$ with the reverse negative, no metric changes significantly: NLI-Score $\Delta{=}{+}0.062$ ($[-0.034,+0.156]$, $p{=}0.44$), hallucination cost $\Delta{=}{-}0.011$ ($p{=}0.21$), omission cost $\Delta{=}{-}0.004$ ($p{=}0.48$), tIoU $\Delta{=}{+}0.021$ ($p{=}0.58$).
Every point estimate is in the hoped-for direction and every interval includes zero.
At $\alpha{=}2.0$ the contrast becomes harmful: hallucination cost $+0.030$ ($p{=}0.031$), omission cost $+0.021$ ($p{=}0.026$), ROUGE-L $-0.016$ ($p{=}0.005$), Full-NLI $-0.098$ ($p{=}0.014$).

\begin{table*}[!t]
\centering
\caption{Paired comparisons on aligned per-video scores ($n{=}150$): mean difference (second system minus first), paired-bootstrap 95\% CI (5{,}000 resamples), two-sided Wilcoxon signed-rank $p$. Arrows give the desired direction. Significant results ($p<0.05$) in \textbf{bold}. Adaptation is significant on every completeness metric, non-significant on \Hcost{}, and significantly \emph{worse} on NLI contradiction mass. Every \tcd{} result at $\alpha{=}0.5$ is non-significant; every significant \tcd{} result at $\alpha{=}2.0$ is a degradation. Remaining variants are in Table~\ref{tab:ho} and Table~\ref{tab:e4}.}
\label{tab:paired}
\scriptsize
\setlength{\tabcolsep}{2.5pt}
\resizebox{\textwidth}{!}{%
\begin{tabular}{@{}l ccc ccc ccc@{}}
\toprule
 & \multicolumn{3}{c}{\textbf{ZS Qwen2.5-VL-7B $\to$ \crashlogic{} greedy}} & \multicolumn{3}{c}{\textbf{Greedy $\to$ \tcd{} $\alpha{=}0.5$ (reverse)}} & \multicolumn{3}{c}{\textbf{Greedy $\to$ \tcd{} $\alpha{=}2.0$ (reverse)}} \\
\cmidrule(lr){2-4}\cmidrule(lr){5-7}\cmidrule(lr){8-10}
\textbf{Metric} & $\Delta$ & 95\% CI & $p$ & $\Delta$ & 95\% CI & $p$ & $\Delta$ & 95\% CI & $p$ \\
\midrule
BLEU-4 (per-video)$\uparrow$ & \textbf{$+$0.108} & $[+.098,+.117]$ & $\mathbf{<10^{-4}}$ & $+$0.000 & $[-.009,+.009]$ & 0.44 & $-$0.009 & $[-.019,+.000]$ & 0.14 \\
ROUGE-L$\uparrow$            & \textbf{$+$0.176} & $[+.163,+.189]$ & $\mathbf{<10^{-4}}$ & $+$0.003 & $[-.005,+.011]$ & 0.22 & \textbf{$-$0.016} & $[-.026,-.006]$ & $\mathbf{0.005}$ \\
tIoU$\uparrow$               & \textbf{$+$0.360} & $[+.297,+.426]$ & $\mathbf{<10^{-4}}$ & $+$0.021 & $[-.044,+.084]$ & 0.58 & $-$0.017 & $[-.083,+.046]$ & 0.58 \\
\Ocost{} (omission)$\downarrow$ & \textbf{$-$0.356} & $[-.382,-.328]$ & $\mathbf{<10^{-4}}$ & $-$0.004 & $[-.017,+.008]$ & 0.48 & \textbf{$+$0.021} & $[+.002,+.040]$ & $\mathbf{0.026}$ \\
\midrule
\Hcost{} (hallucination)$\downarrow$ & $-$0.005 & $[-.051,+.041]$ & 0.81 & $-$0.011 & $[-.027,+.005]$ & 0.21 & \textbf{$+$0.030} & $[+.005,+.057]$ & $\mathbf{0.031}$ \\
NLI $P_e$$\uparrow$          & \textbf{$+$0.084} & $[+.019,+.148]$ & $\mathbf{0.007}$ & $+$0.026 & $[-.029,+.080]$ & 0.24 & $-$0.005 & $[-.063,+.055]$ & 0.51 \\
NLI $P_c$$\downarrow$        & \textbf{$+$0.137} & $[+.069,+.203]$ & $\mathbf{<10^{-4}}$ & $-$0.036 & $[-.094,+.022]$ & 0.94 & $-$0.027 & $[-.089,+.034]$ & 0.58 \\
NLI-Score$\uparrow$          & $-$0.052 & $[-.162,+.055]$ & 0.41 & $+$0.062 & $[-.034,+.156]$ & 0.44 & $+$0.021 & $[-.075,+.119]$ & 0.91 \\
Full-NLI-Score$\uparrow$     & \textbf{$+$0.170} & $[+.077,+.269]$ & $\mathbf{0.0005}$ & $-$0.012 & $[-.104,+.081]$ & 0.77 & \textbf{$-$0.098} & $[-.186,-.013]$ & $\mathbf{0.014}$ \\
\bottomrule
\end{tabular}%
}
\end{table*}


A methodological note: an unpaired normal-approximation interval with a pooled standard deviation of $0.25$ would call the $\alpha{=}0.5$ NLI gain significant ($z\approx2.1$).
The actual per-video standard deviation of NLI-Score is $\approx0.55$, and the paired test does not support the claim~\cite{dror2018hitchhiker}.
We therefore treat mild \tcd{} as \emph{safe but ineffective} on this task and strong \tcd{} as \emph{harmful}, and we spend the next section explaining why.
